\documentclass[sigconf,nonacm]{acmart}
\settopmatter{printacmref=false,printfolios=true}
\usepackage{xurl}
\setlength{\emergencystretch}{2em}
\title{Replaying Freeway Estimates After Late Detector Corrections}
\subtitle{ACM SIGMOD 2027 | Demonstration draft}
\author{Arjun Vijay Prakash}
\affiliation{\institution{Independent Researcher}\city{ }\country{ }}

\begin{document}
\begin{abstract}
A late correction to a freeway detector record can change several estimated windows. Rebuilding
an entire replay repeats work, and changing only the corrected window can preserve estimates
that depended on the old value. CREST stores observations as immutable, content-addressed
Parquet partitions. Each estimate records its input partitions and the topology and estimator
settings used to compute it. The demonstration lets attendees mask a detector reading, compare
the old and revised manifests, and verify the new output against full recomputation. Two
retrospective Caltrans PeMS replays recomputed 4 of 16 and 4 of 17 windows, with exact output
parity in both traces.
\end{abstract}

\maketitle

\section*{The system}
An input manifest maps event timestamps to immutable Parquet files. Each file records one time
partition of detector speeds, with missing values represented as nulls. A correction creates a
new input manifest and changes the partition digest for the affected timestamp. Earlier
manifests remain usable, so a reviewer can recover the evidence behind a previous estimate.

A result cache key covers the output timestamp, all input partition digests in its dependency
window, the directed topology, numerical settings, and estimator and solver versions. The
default estimator uses the current partition and up to three preceding partitions. For an
estimator configured with lookahead, dependencies also include the requested future partitions.
The replay checks every output key and rebuilds affected windows. As a result, numerical work
can stay local even though manifest traversal still scales with the trace length. The current
system supports bounded-window estimates with no state carried beyond the declared dependency
window.

\section*{What the attendee does}
First, load the original manifest and choose a timestamp and station. Inspect the detector
speed, estimate, solver status and dependency digests. Next, remove that reading using the
correction command and inspect the revised manifest. Compare the affected outputs and follow
their dependency digests to the changed source partition. Apply the same edit again to reproduce
the corrected manifest. Finally, rebuild the complete corrected trace and compare it with the
incremental output, including the estimates and their provenance.

The demonstration uses a clean CPU installation and the bundled synthetic junction scenario when
the original measured snapshots are unavailable. The command-line interaction runs the estimator
after a correction. The measured traces remain local pending source redistribution review.

\section*{Evidence and scope}
On I-880 southbound, the correction recomputed 4 of 16 outputs and preserved 12. On I-10
westbound it recomputed 4 of 17 and preserved 13. All revised outputs matched a full replay. In
the 72-window synthetic junction example, 4 were recomputed and 68 were preserved, again with
exact parity. A regression fixture repeats a correction after later partitions have been
ingested and checks idempotence and configured lookahead. Bulk ingestion expects a sorted,
unique time grid.

These traces establish the invalidation behavior and demonstrate the attendee interaction. The
current evidence describes work avoided on short, single-corridor traces. A larger trace would
be needed to characterize end-to-end speedup, including partition I/O and cache traversal.

\section*{Reproduction}
Run \texttt{python -m src.research.roadmap fixture} to generate a synthetic replay and its
receipt. The measured workflow, correction syntax, and packaged development receipt are linked
from the implementation guide.

\section*{References}
[1] Caltrans. Performance Measurement System. \url{https://pems.dot.ca.gov/}

[2] CREST implementation guide and development receipts.
\url{https://github.com/ArjunCodess/CREST/blob/main/docs/crest-implementation.md}
\end{document}
